\subsection{Uniform truncated Gaussian approximation}
\label{sec:gaussian_uniform_ground}

% Source: polynomial-PEPS manuscript (September 24, 2026), lem:reset,
% pinned peps-02-information-adc7f124.tex, lines 425--452, at
% openai/math@adc7f1241b42e322a6451854ab7e4b4c146bf78a.
% This is the approximation step only, with arbitrary complex overlap.
% No spatial locality or full reset result is claimed.
% Production, consumers, and axiom reports checked at 3702df8.
% This leaf is not added to a shared router in the present change.

\begin{theorem}[Uniform truncated Gaussian approximation]
    \label{thm:gaussian_uniform_ground}
    \lean{GaussianFilter.exists_uniform_truncated_ground_estimate}
    \leanok
    Let $\Delta>0$ and $r\ge0$ be real. There exist $A>0$ and $T_0>0$,
    depending only on $\Delta,r$, with the following property. For every
    real $b\ge0$ and $L\ge1$, put $B_*=1+b+\log L$, $h=AB_*$, and
    $T=T_0B_*$. Then $h>0$ and $T\ge0$.

    In any finite-dimensional complex Hilbert space, let $H,H'$ be
    self-adjoint operators and let $W$ be a contraction. Suppose that
    unit vectors $\Psi,\Phi$ and a real number $E_0$ satisfy
    \begin{align}
        H\Psi&=E_0\Psi,& H'\Phi&=E_0\Phi,
        \label{eq:gaussian_uniform_ground_vectors}
    \end{align}
    \begin{align}
        H-E_0\Id&\ge\Delta(\Id-\ketbra{\Psi}{\Psi}),&
        H'-E_0\Id&\ge\Delta(\Id-\ketbra{\Phi}{\Phi}).
        \label{eq:gaussian_uniform_ground_gaps}
    \end{align}
    Define $z=\braket{\Phi}{W\Psi}$ and the unrenormalized integral
    \begin{align}
        M_T&=\int_{-T}^{T}(2\pi h)^{-1/2}e^{-t^2/(2h)}
            e^{itH'}W e^{-itH}\,dt.
        \label{eq:gaussian_uniform_ground_integral}
    \end{align}
    Then $\|M_T\|\le1$, and
    \begin{align}
        \|M_T\Psi-z\Phi\|&\le3e^{-10b}L^{-r},
        \label{eq:gaussian_uniform_ground_forward}
    \end{align}
    \begin{align}
        \|M_T^\dagger\Phi-\overline z\Psi\|&\le3e^{-10b}L^{-r}.
        \label{eq:gaussian_uniform_ground_reverse}
    \end{align}
    This is the uniform Gaussian approximation in the reset argument
    of~\cite{OpenAI2026PolynomialPEPS}.
\end{theorem}
\begin{proof}\leanok
    \uses{thm:gaussian_parameters, thm:gaussian_filter_contraction,
        thm:gaussian_truncated_two_sided}
    Choose $A,T_0$ by Theorem~\ref{thm:gaussian_parameters}, before
    $b,L$, the Hilbert space, and the operators and vectors. The positive
    Gaussian density has total mass one, and each operator
    $e^{itH'}W e^{-itH}$ has norm at most $\|W\|\le1$.
    Restricting its integral to $[-T,T]$ therefore gives $\|M_T\|\le1$.

    Theorem~\ref{thm:gaussian_truncated_two_sided}
    bounds each residual in (\ref{eq:gaussian_uniform_ground_forward}) and
    (\ref{eq:gaussian_uniform_ground_reverse}) by
    \begin{align}
        \bigl(e^{-h\Delta^2/2}+2e^{-T^2/(2h)}\bigr)\|W\|
        &\le \bigl(1+2\bigr)e^{-10b}L^{-r}
         =3e^{-10b}L^{-r}.
        \label{eq:gaussian_uniform_ground_error_sum}
    \end{align}
    The last inequality uses
    (\ref{eq:gaussian_parameter_error}) and
    (\ref{eq:gaussian_parameter_truncation}). Conjugate symmetry gives
    $\braket{\Psi}{W^\dagger\Phi}=\overline z$ in the adjoint estimate.
\end{proof}
